\subsection{Motivation}

Our primary motivation to write this paper comes from two questions, which are broadly related under the microcosm principle. The first one concerns the dualities of categorical structures. Consider the $\op_1$ duality on $\nCat{1}$, expressed as a $2$-functor $\op_1 \colon \op_2 (\nCat{1}) \rightarrow \nCat{1}$, which satisfy the relation $\op_1 \circ \op_2 (\op_1) =1$. Clearly, this is not an involution in the traditional sense of satisfying $\op_1 \circ \op_1 = 1$, since this composition is simply not defined. Rather, it is an involution relative to a higher involution $\op_2$. This motivates the definition of a relative involution.

Now consider the duality $\op_1 \colon \op_2(\MonCat_{c}) \rightarrow \MonCat_{l}$. Here, the domain consists of colax monoidal functors, whereas the codomain consists of lax monoidal functors. This is again a duality, but it is no longer directly expressible even using relative involution as before. Instead, this motivates a further generalization to relative inverses. 

As with many algebraic structures, specifying the complete coherence data for these definitions becomes increasingly difficult in higher dimensions. Nevertheless, it is tractable and useful enough to define these for involutive bicategories and biadjunction of bicategories. In particular, this framework is sufficiently general to express the dualities of monoidal structures. Furthermore, it is also natural to consider module structures in the presence of monoidal structures. Thus, one of the main goals of this paper is to extend this framework onto virtual double categories and the corresponding module construction.  

The second question concerns slicing and its interaction with these dualities. In particular, we wish to describe a general construction of enriched slice and coslice categories, especially under the consideration of model structures. Although there are some accounts of enriched slice constructions, it is useful to have a systematic treatment that encompasses both slice and coslice constructions. Furthermore, the slice and coslice category construction is related by the $\op_1$ duality. This suggests that the enriched coslice case should be a dual of the enriched slice case in some appropriate sense. Establishing such a construction is the second main goal of this paper.


\subsection{Main results and applications}

The results can be divided into three major parts. The first part establishes a desired framework of relative involutions and relative inverses. The second part applies this framework to characterize dualities of constructions on virtual double categories. The third part discuss the slicing of module structures and establish a general construction of enriched slice and coslice model categories.

Recall that one can define a $1$-category of relative involutions, called $\nInv{0}(\cC,T,\alpha)$, inside an arbitrary lax involutive category $(\cC,T,\alpha)$ \cite{yauinvcat2020} (or see \cref{sub:1diminv}). This is an example of the microcosm principle. Motivated by this fact, we show a similar construction for adjunctions by defining relative inverses. The main result in this part is to generalize this framework one dimension higher for weak $2$-involutive bicategories and weak $2$-inverses of bicategories. This can be summarized by the following proposition.

\begin{proposition*}
	The microcosm principle holds for weak $2$-involutive bicategories and weak $2$-inverses of bicategories.
\end{proposition*}

In particular, this provides suitable definitions of relative weak $1$-inverses, relative weak $1$-involutions, and their $1$-morphisms. Note that expressing the dualities of strict $n$-categories only require relative involution or relative inverse (see \cref{prop:opdualities}). This is in contrast to the monoidal structures, which also require the $1$-morphisms. The main result here is that the $2$-functors $\Mod \colon \vDblCat \rightarrow \vDblCat$ and $\CoMod \colon \vcoDblCat \rightarrow \vcoDblCat$ are compatible with the relevant dualities.

\begin{theorem*} The following extensions hold:
	\begin{enumerate}
		\item $\Mod$ extends to a $1$-morphism of involution, respecting the horizontal duality of virtual double categories. 
		
		\item $\Mod$ and $\CoMod$ together extend to a $1$-morphism of inverses, respecting the vertical duality of virtual double categories.
	\end{enumerate}
\end{theorem*}

In the final part, we first briefly review the slice constructions for module categories. Recall that standard definitions of an enriched model category assumes the existence of a two-variable closure adjunction relating tensoring, powering, and enrichement \cite{hovey2007modelcat} (or see \cref{sub:slicecatmod}). Hence, we focus primarily on the behavior of this two-variable adjunction under slicing. We show that it can be sliced in two dual ways, one by slicing the tensoring and the other by slicing the powering. The main result is that these slicing also respect the model structures.

\begin{theorem*}
	Let $\ucC$ be a $\ucV$-enriched model categories. Let $M\in \cV$ and $Z \in \cC$ be suitable monoid and module object, respectively. Then, both $Z/\ucC$ and $\ucC/Z$ are $\ucV/M$-enriched model categories.
\end{theorem*}

\subsection{Overview}

\Cref{sec:prelim} discusses the dualities of strict $n$-categories as a motivating example of nested algebraic structures. To formulate this more precisely, it recall the microcosm principle and define the sliced microcosm principle. A relevant example of these principles is then provided in terms of pseudomonoid and pseudomodule objects.

\Cref{sec:ginvs} begins by briefly reviewing the theory of involutions and relative involutions. Correspondingly, it defines the similar theory of inverses and relative inverses. Then, the analogues of these structures are described in one dimension higher. 

\Cref{sec:dualvdcs} starts with reformulating the dualities of monoidal categories within the framework of relative inverses and involutions. This framework is then used to describe the corresponding dualities of internal monoid structures. Subsequently,  the dualities are generalized to the level of virtual double categories, encoding the action on the module data.

\Cref{sec:slicecatmod} applies the sliced microcosm principle to monoidal and module structures. First, it reviews enrichement through tensoring and states the corresponding slice construction. The results are then dualized for the case of cotensoring. Finally, these constructions are extended to the setting of enriched model categories. 